\documentclass{article}

% -------------------------------------------------
% Page layout
% -------------------------------------------------
\usepackage[
    a4paper,
    margin=1.5cm
]{geometry}

% -------------------------------------------------
% Typography
% -------------------------------------------------
\usepackage[T1]{fontenc}
\usepackage{newtxtext}
\usepackage{newtxmath}
\usepackage{microtype}

% -------------------------------------------------
% Mathematics
% -------------------------------------------------
\usepackage{amsmath}


% -------------------------------------------------
% Figures and tables
% -------------------------------------------------
\usepackage{graphicx}
\usepackage{booktabs}
\usepackage{array}
\usepackage[table]{xcolor}
\usepackage{tikz}
\usetikzlibrary{fit,positioning,arrows.meta,calc,shadows.blur,shapes.geometric,backgrounds}
\usepackage{fontawesome5}
\usepackage{caption}
\usepackage{capt-of}

% -------------------------------------------------
% References
% -------------------------------------------------
\usepackage[numbers]{natbib}
\usepackage{hyperref}

% -------------------------------------------------
% Table styling
% -------------------------------------------------
\definecolor{toscarow}{RGB}{252,241,241}

\captionsetup[table]{
    font=small,
    labelfont=normalfont,
    textfont=normalfont,
    justification=justified,
    singlelinecheck=false,
    skip=5pt
}

% Regular result: mean ± standard deviation
\newcommand{\result}[2]{%
    \ensuremath{#1\,{\scriptstyle \pm\,#2}}%
}

% Best result: bold mean, regular smaller standard deviation
\newcommand{\bestresult}[2]{%
    \ensuremath{\mathbf{#1}\,{\scriptstyle \pm\,#2}}%
}

\title{Adini Gate Koydum}
\author{Elif Ceren Gok}
\date{July 2026}

\begin{document}

\maketitle

\section{Introduction}

Introduction

\section{Related Work}
\label{sec:related-work}

\paragraph{CIL with Randomly Initialized Models.}
Before foundation models became the default backbone choice, most
class-incremental learning (CIL) research trained convolutional or residual
networks \citep{he2016deep} from random weights, falling into four broad
directions. Regularization-based approaches restrict how much important
parameters may move between tasks, either penalizing weight change directly
\citep{kirkpatrick2017overcoming,zenke2017continual} or keeping predictions
close to an earlier snapshot of the model
\citep{li2018learning,aljundi2018memory}. Replay-based approaches keep a
small window into the past, either storing a subset of old exemplars
\citep{rebuffi2017icarl,lopezpaz2017gradient,bang2021rainbow,prabhu2020gdumb,wu2019large,zhao2020maintaining,liu2021rmm,arani2022learning,sarfraz2023error}
or training a generator that reproduces old-class samples on demand
\citep{shin2017continual,he2018exemplarsupported,hu2019overcoming,petit2023fetril,mocanu2016online}.
A fourth direction reshapes the network itself, growing fresh capacity for
every task
\citep{aljundi2017expert,yan2021der,wang2022foster,wang2023task,zhou2023model}
or pruning a shared network into task-specific sub-networks
\citep{mallya2018packnet,wortsman2020supermasks,kang2022forgetfree,kang2023soft,golkar2019continual,wang2022sparcl,yildirim2024continual,yildirim2025selfregulated}.

\paragraph{CIL with Pre-trained Models.}
The picture changes considerably once a strong pre-trained backbone,
typically a Vision Transformer \citep{dosovitskiy2021image}, is available
from the start, since it already generalizes well before any task-specific
training happens at all \citep{ramasesh2022effect,mehta2023empirical}. The
open question shifts from \emph{how do we prevent forgetting while learning
from scratch} to \emph{how little can we touch the model without losing
that generalization} \citep{wu2022classincremental}, and the resulting work
again falls into a few directions. Prompting methods leave the backbone
fully untouched and learn a small set of input-level tokens instead, first
routed per instance through a learned key-query pool \citep{wang2022learning},
then refined with task-shared/task-specific splits
\citep{wang2022dualprompt}, attention-based reweighting
\citep{smith2023coda}, and feature-space replay imitation
\citep{wang2023hierarchical}. A second direction fits only a classifier on
top of frozen features, either through simple class prototypes
\citep{zhou2024revisiting} or, as in RanPAC \citep{mcdonnell2023ranpac},
through a closed-form ridge classifier on randomly projected features,
building on the first-session training stage of \citet{panos2023first}. A
third direction inserts adapters into the backbone directly, either
layer-wise and accumulated across tasks
\citep{zhou2024expandable,sun2025mos}, or, as in TOSCA
\citep{yildirim2026unlocking}, as a single sparse adapter-calibrator pair
(LuCA) applied only to the \texttt{[CLS]} token right before the
classifier, with a fresh pair trained per task and every stored module
scored at inference to pick the lowest-entropy prediction. Two of these
methods, RanPAC and TOSCA, are the direct starting point for this work, as
discussed in Section~\ref{sec:method}.

\section{Background}
\label{sec:background}

This section lays out the CIL setting formally and then walks through the
three main families of post-training strategy that build on top of a
frozen foundation model, closing with the specific limitation in each
family that motivates the design in Section~\ref{sec:method}.

\subsection{Class-Incremental Learning (CIL)}
\label{sec:background-cil}

In CIL, a model is exposed to a sequence of $T$ disjoint tasks
$\{\mathcal{T}_1, \ldots, \mathcal{T}_T\}$, trained on one task at a time.
Task $t$ comes with its own labeled dataset
$\mathcal{D}_t=\{(x_i,y_i)\}_{i=1}^{n_t}$, drawn from classes $Y_t$
disjoint from every earlier task's, $Y_t \cap Y_{t'}=\emptyset$ for $t \neq
t'$. Once training on $\mathcal{D}_t$ starts, data from earlier tasks is no
longer accessible, which is what makes the setting hard: the model must
keep working on all classes seen so far, $\mathcal{Y}=Y_1 \cup \cdots \cup
Y_t$, using only what it retained internally.

Following the FM-based CIL literature
\citep{wang2022learning,wang2022dualprompt,smith2023coda,zhou2024revisiting,zhou2024expandable,sun2025mos},
we write the model as $f(\mathbf{x}) = W^\top \Phi(\mathbf{x})$, splitting
it into a feature extractor $\Phi: \mathcal{X} \to \mathbb{R}^d$ and a
classifier $W \in \mathbb{R}^{d \times |\mathcal{Y}|}$. For a ViT
\citep{dosovitskiy2021image}, $\Phi$ patchifies the input into a token
sequence, prepends a learnable \texttt{[CLS]} token, and passes everything
through self-attention and feed-forward blocks; by convention, only the
final \texttt{[CLS]} representation is kept as $\Phi(\mathbf{x})$, since it
is trained to aggregate information from the whole image. Throughout, we
assume the replay-free setting
\citep{wang2022learning,wang2022dualprompt,smith2023coda,zhou2024revisiting,zhou2024expandable}
--- the setting TOSCA, RanPAC, and [METHOD] all operate in --- so no
exemplars from earlier tasks are stored, and the task identity of a test
sample is never revealed at inference: the model must infer both
\emph{what} the object is and, implicitly, \emph{which task it came from}.

\subsection{Overview of Post-Training in CIL}
\label{sec:background-posttraining}

Once the backbone is frozen, essentially every method in this space reduces
to the question of what small piece to attach on top of $\Phi(\mathbf{x})$,
and how to train it without disturbing $\Phi$ itself. We group these
strategies into three families.

\paragraph{Learning Prototypical Classifiers.}
The simplest strategy skips learned parameters almost entirely: for every
class $y$ seen in task $t$, a prototype averages the frozen features of
that class's $n_y$ training examples,
\begin{equation}
  \mathbf{p}_y = \frac{1}{n_y}\sum_{i:\,y_i=y} \Phi(\mathbf{x}_i),
  \label{eq:prototype}
\end{equation}
and a test sample is classified by whichever prototype it lies closest to.
RanPAC \citep{mcdonnell2023ranpac} keeps this recipe but strengthens the
distance computation: features are first passed through a fixed random
projection with a nonlinearity, $\phi(\mathbf{x}) =
\mathrm{ReLU}(\Phi(\mathbf{x})^\top P)$ with $P \in \mathbb{R}^{d \times
M}$, $M \gg d$, and the classifier weights are obtained in closed form as
\begin{equation}
  W^\star = (G + \lambda I)^{-1} C,
  \label{eq:ridge}
\end{equation}
where $G$ and $C$ are Gram and class-sum matrices accumulated over
$\phi(\mathbf{x})$ rather than solved by gradient descent. Together with
the first-session training stage described above, this lets RanPAC skip
training a classifier from scratch for every incoming task --- it only
ever updates $G$ and $C$. The trade-off is that once the first task's
adaptation is frozen, the feature space no longer moves, so the only thing
that can still adjust to new classes is the accumulation of statistics on
top of it; genuinely task-specific structure arriving after task one has
nowhere to be represented.

\paragraph{Learning Prompts.}
Rather than modifying the classifier, this family modifies what the
backbone sees as input: a pool of learnable prompt vectors
$\mathcal{P}=\{P_1,\dots,P_M\}$, each paired with a trainable key, is
inserted into the embedding layer or across several transformer blocks,
and an instance is routed to the prompts whose keys are closest to its
query embedding under the \emph{frozen} backbone
\citep{wang2022learning,wang2022dualprompt,smith2023coda,wang2023hierarchical}.
Because prompts are shared rather than one-per-task, this keeps parameter
count small and the backbone untouched, a real advantage for stability.
The difficulty shows up at retrieval time: as tasks accumulate, prompts
from different tasks can end up describing similar-looking classes, and
since the key space is fixed from the start, it has no way to later
separate prompts that turn out to conflict --- exactly the kind of
interference that leads to forgetting in this family.

\paragraph{Learning Adapters.}
The third family edits the backbone directly, inserting small bottleneck
modules into its layers. For an intermediate representation $\mathbf{z}$, a
task-specific adapter set $\mathcal{A}_t = \{A_1, \ldots, A_N\}$ (one module
per transformer layer, in the layer-wise case) applies
\begin{equation}
  A(\mathbf{z}) = \mathbf{z} + \psi(\mathbf{z}W_{down})W_{up}, \qquad
  W_{down}\in\mathbb{R}^{d\times r},\ W_{up}\in\mathbb{R}^{r\times d},
  \label{eq:adapter}
\end{equation}
with $\psi$ a nonlinearity and $r \ll d$
\citep{zhou2024revisiting,zhou2024expandable,sun2025mos}. Placing one
module in every layer, for every task, is effective but expensive:
parameter count grows as $\mathcal{O}(T\!\cdot\! N\!\cdot\! dr)$, and small
shifts introduced at every layer tend to compound by the time they reach
the classifier. TOSCA \citep{yildirim2026unlocking} avoids layer-wise
placement entirely, applying a single adapter-calibrator pair only to the
final \texttt{[CLS]} representation,
\begin{equation}
  \Phi(\mathbf{x})' = C(A(\Phi(\mathbf{x}))),
  \label{eq:tosca-clsonly}
\end{equation}
which keeps the parameter count fixed at $\mathcal{O}(dr)$ regardless of
depth, and trains it with an added $\ell_1$ penalty that pushes different
tasks' modules toward disjoint, sparse sets of active weights. The cost
shows up at inference instead: with no task label available, TOSCA must
run every stored module on a given input and keep the prediction with the
lowest entropy, so scoring one test sample takes work proportional to the
number of tasks seen so far.

\section{Methodology}
\label{sec:method}

\input{resources/methodology_figure_tikz.tex}


Our method combines the task-specific adaptation mechanism of TOSCA with a
closed-form ridge-based routing and classification pipeline inspired by
RanPAC. The guiding idea is simple: instead of evaluating every stored
task-specific module at test time, we first route an input to the most
likely task through a global classifier trained in a shared frozen feature
space, and only then invoke the corresponding task expert. In this way, we
retain TOSCA's plasticity while removing its linear-in-\(T\) inference
cost.



\subsection{Overview}
\label{subsec:method-overview}

Let \(f_{\theta}\) denote a ViT-B/16 backbone pretrained on ImageNet-21K,
and let \(\Phi(\mathbf{x}) \in \mathbb{R}^{d}\) be the final
\texttt{[CLS]} token representation extracted from input \(\mathbf{x}\).
At task \(t\), we learn a lightweight task-specific LuCA module
\(E_t\) on top of \(\Phi(\mathbf{x})\), while keeping the vast majority of
the backbone fixed. Each task therefore contributes two components:



\begin{enumerate}
    \item a \emph{task-specific TOSCA expert} \(E_t\), which refines the
    frozen \texttt{[CLS]} representation for that task; and
    \item a \emph{task-local ridge head}, solved in closed form on top of the
    adapted features of that task.
\end{enumerate}

In parallel, we maintain a \emph{global ridge router} trained on frozen ViT
features shared by all tasks. At inference time, this global router predicts
the most likely task for a sample, and the corresponding expert is then used
for final classification.

\subsection{Task-Specific Token Adaptation}
\label{subsec:method-tosca}

Following TOSCA, each task is assigned a sparse LuCA module applied only to
the final \texttt{[CLS]} token. Given the frozen token feature
\(\Phi(\mathbf{x})\), the adapter branch produces a task-dependent residual
update,
\begin{equation}
    h_a = W^{a}_{up}\,\sigma\!\left(W^{a}_{down}\,\Phi(\mathbf{x})\right),
\end{equation}
where \(W^{a}_{down} \in \mathbb{R}^{d/r_m \times d}\) and
\(W^{a}_{up} \in \mathbb{R}^{d \times d/r_m}\), and \(r_m\) is the adapter
bottleneck ratio.

The calibrator branch computes a channel-wise gating vector,
\begin{equation}
    g = \mathrm{Sigmoid}\!\left(
    W^{c}_{up}\,\delta\!\left(W^{c}_{down}\,\Phi(\mathbf{x})\right)
    \right),
\end{equation}
where \(W^{c}_{down} \in \mathbb{R}^{d/r_s \times d}\),
\(W^{c}_{up} \in \mathbb{R}^{d \times d/r_s}\), and \(r_s\) is the
calibration bottleneck ratio. The final adapted representation is
\begin{equation}
    \tilde{\Phi}_t(\mathbf{x}) =
    \left(\Phi(\mathbf{x}) + \mathrm{LN}(h_a)\right) \odot g.
\end{equation}

This design keeps the adaptation narrowly focused on the last token instead
of altering all transformer layers. As a result, the backbone preserves the
general semantic structure learned during pretraining, while each task still
receives a compact task-specific correction.

\subsection{Incremental Training Objective}
\label{subsec:method-training}

At incremental step \(t\), only the current task data \(\mathcal{D}_t\) is
available. We freeze the entire backbone except for the TOSCA parameters and
the classifier parameters associated with the current training stage. During
the first task, the adapter parameters inside the ViT are additionally
allowed to update once for initial specialization; afterwards they remain
frozen for the rest of the sequence.

Training for task \(t\) minimizes a cross-entropy objective over the newly
introduced classes,
\begin{equation}
    \mathcal{L}_{\mathrm{ce}} =
    - \sum_{(\mathbf{x},y)\in\mathcal{D}_t}
    \log p(y \mid \mathbf{x}),
\end{equation}
augmented with an \(\ell_1\) penalty over the TOSCA parameters,
\begin{equation}
    \mathcal{L} =
    \mathcal{L}_{\mathrm{ce}} + \lambda \|\Theta_{\mathrm{TOSCA}}\|_1,
\end{equation}
where \(\Theta_{\mathrm{TOSCA}}\) denotes the parameters of the active
task-specific LuCA module and \(\lambda\) is a sparsity coefficient. This
regularization encourages each task expert to remain compact and discourages
overlap between task-specific adaptations.

After optimizing task \(t\), the resulting TOSCA parameters are stored as an
immutable expert \(E_t\). The current task's classifier is then replaced by
a ridge head computed from the adapted task features, as described next.

\subsection{Shared Random Projection and Ridge Features}
\label{subsec:method-ridge-features}

To make closed-form ridge classification effective, we follow the
RanPAC-style feature lifting strategy. Let \(z = \Phi(\mathbf{x})\) denote
the frozen ViT feature before task-specific adaptation. We first optionally
normalize \(z\), then project it into a higher-dimensional random feature
space:
\begin{equation}
    \phi(\mathbf{x}) = \mathrm{ReLU}(z^\top P),
\end{equation}
where \(P \in \mathbb{R}^{d \times M}\) is a fixed random projection matrix
sampled once at the beginning of training and shared across all tasks. The
projection dimension \(M\) is significantly larger than \(d\), allowing a
linear ridge head in this lifted space to model richer class boundaries
without gradient-based optimization.

Importantly, the same projection matrix is used for all tasks and for the
global router. This creates a common feature space in which both routing and
task-local classification can be expressed through closed-form solvers.

\subsection{Task-Local Ridge Heads}
\label{subsec:method-local-ridge}

For every task \(t\), we collect the adapted features
\(\tilde{\Phi}_t(\mathbf{x})\) of current-task training samples, map them
through the shared random projection, and obtain
\(\phi_t(\mathbf{x}) \in \mathbb{R}^{M}\). Let \(\Phi_t\) be the matrix that
stacks these projected features, and let \(Y_t\) be the task-local one-hot
label matrix over the classes of task \(t\). The task-specific ridge head is
then given by
\begin{equation}
    W_t =
    \left(\Phi_t^\top \Phi_t + \lambda_r I\right)^{-1}\Phi_t^\top Y_t,
\end{equation}
where \(\lambda_r\) is the ridge regularization coefficient.

Rather than retraining a parametric classifier with gradient descent, we only
need to accumulate the corresponding Gram and class-sum statistics,
\begin{equation}
    G_t = \Phi_t^\top \Phi_t,
    \qquad
    C_t = \Phi_t^\top Y_t,
\end{equation}
from which the ridge solution can be re-solved exactly whenever needed. This
makes the classifier fitting stage fast, stable, and independent of stored
exemplars.

\subsection{Global Ridge Router}
\label{subsec:method-global-router}

The central mechanism that removes TOSCA's inference bottleneck is a global
ridge router trained in the \emph{frozen} ViT feature space. For every
sample observed so far, we extract its frozen feature \(z=\Phi(\mathbf{x})\),
apply the same random projection to obtain \(\phi(\mathbf{x})\), and
accumulate global statistics over all classes seen up to task \(t\):
\begin{equation}
    G^{\mathrm{glob}} = \Phi_{\mathrm{glob}}^\top \Phi_{\mathrm{glob}},
    \qquad
    C^{\mathrm{glob}} = \Phi_{\mathrm{glob}}^\top Y_{\mathrm{glob}}.
\end{equation}
The corresponding closed-form global ridge classifier is
\begin{equation}
    W^{\mathrm{glob}} =
    \left(G^{\mathrm{glob}} + \lambda_r I\right)^{-1} C^{\mathrm{glob}}.
\end{equation}

Given a test input \(\mathbf{x}\), the global router first produces
class-level scores over the union of all seen classes,
\begin{equation}
    s^{\mathrm{glob}}(\mathbf{x}) = \phi(\mathbf{x})^\top W^{\mathrm{glob}}.
\end{equation}
These class scores are then converted into task scores by taking the maximum
class score inside each task block:
\begin{equation}
    s_t(\mathbf{x}) =
    \max_{c \in \mathcal{C}_t} s^{\mathrm{glob}}_c(\mathbf{x}).
\end{equation}
The routed task is selected as
\begin{equation}
    \hat{t} = \arg\max_t s_t(\mathbf{x}).
\end{equation}

This routing mechanism is exemplar-free and closed-form: no task labels are
needed at inference, and no additional gradient-based router training is
required after the backbone features are extracted.

\subsection{Inference}
\label{subsec:method-inference}

Inference proceeds in two stages. First, the global ridge router predicts the
most likely task \(\hat{t}\) from frozen ViT features. Second, only the
corresponding stored TOSCA expert \(E_{\hat{t}}\) is loaded and applied to
the same sample, yielding an adapted task-specific representation that is
scored by the corresponding ridge head \(W_{\hat{t}}\). The final prediction
is therefore
\begin{equation}
    \hat{y} =
    \arg\max_{c \in \mathcal{C}_{\hat{t}}}
    \phi_{\hat{t}}(\mathbf{x})^\top W_{\hat{t}}.
\end{equation}

This design avoids TOSCA's original test-time strategy of evaluating every
stored expert and choosing the one with the lowest entropy. Instead, only a
single expert is invoked per sample after routing, which makes the test-time
cost effectively constant with respect to the number of tasks, aside from the
single global ridge pass.

\subsection{Why This Combination Works}
\label{subsec:method-discussion}

The architecture is built around a division of labor between three parts:

\begin{itemize}
    \item the \emph{frozen ViT backbone}, which preserves general
    visual-semantic knowledge and supplies stable task-agnostic features;
    \item the \emph{task-specific TOSCA experts}, which provide lightweight
    plasticity only where adaptation is needed; and
    \item the \emph{ridge-based routing and classification heads}, which turn
    both routing and classification into stable closed-form problems rather
    than gradient-trained modules.
\end{itemize}

Compared to RanPAC, the method restores ongoing task-specific adaptation by
retaining a separate lightweight expert for each task. Compared to TOSCA, it
eliminates the need to score all experts at inference. In this sense,
[METHOD] combines TOSCA's plasticity with RanPAC's closed-form efficiency,
resulting in a replay-free, parameter-efficient, and test-time scalable
foundation-model-based CIL framework.

\section{Experiments}
\label{sec:experiments}

We evaluate on six benchmarks following the standard FM-based CIL protocol
\citep{zhou2024revisiting,zhou2024expandable}, using ViT-B/16 pretrained on
ImageNet-21K as the frozen backbone throughout. Table~\ref{tab:main-results}
reports accuracy and Table~\ref{tab:efficiency} reports computational cost.

\begin{table*}[htbp]
\centering

\caption{
Average accuracy ($\bar{A}$) and last accuracy ($A_T$) on six datasets
using ViT-B/16 pretrained on ImageNet-21K.
Results are reported as mean $\pm$ standard deviation over five seeds.
The best performance in each column is shown in bold.
}

\label{tab:main-results}

\small
\setlength{\tabcolsep}{5pt}
\renewcommand{\arraystretch}{1.08}

\resizebox{\textwidth}{!}{%
\begin{tabular}{
    l
    *{12}{c}
}
\toprule

& \multicolumn{2}{c}{CIFAR B0 Inc5}
& \multicolumn{2}{c}{CUB B0 Inc10}
& \multicolumn{2}{c}{IN-R B0 Inc20}
& \multicolumn{2}{c}{IN-A B0 Inc20}
& \multicolumn{2}{c}{OmniBench B0 Inc30}
& \multicolumn{2}{c}{VTAB B0 Inc10}
\\[-1pt]

Method
& $\bar{A}$ & $A_T$
& $\bar{A}$ & $A_T$
& $\bar{A}$ & $A_T$
& $\bar{A}$ & $A_T$
& $\bar{A}$ & $A_T$
& $\bar{A}$ & $A_T$
\\

\midrule

SimpleCIL
& \result{86.48}{0.72} & \result{81.28}{0.00}
& \result{91.62}{1.12} & \result{86.77}{0.00}
& \result{61.40}{0.72} & \result{55.11}{0.00}
& \result{58.67}{0.90} & \result{48.39}{0.00}
& \result{79.68}{1.33} & \result{73.30}{0.00}
& \result{90.65}{0.98} & \result{84.45}{0.00}
\\

RanPAC
& \result{93.43}{0.63} & \result{90.02}{0.27}
& \result{93.09}{0.61} & \result{89.27}{0.49}
& \result{83.87}{0.52} & \result{78.86}{0.17}
& \result{63.47}{1.45} & \result{54.73}{0.73}
& \result{85.84}{0.87} & \result{79.44}{0.35}
& \result{96.12}{0.56} & \result{92.68}{0.29}
\\

L2P
& \result{84.55}{1.74} & \result{79.75}{1.15}
& \result{73.96}{2.00} & \result{62.38}{0.42}
& \result{75.60}{0.59} & \result{70.93}{0.20}
& \result{48.21}{0.83} & \result{41.08}{0.96}
& \result{73.74}{1.74} & \result{64.65}{1.00}
& \result{75.38}{5.53} & \result{59.37}{6.79}
\\

DualPrompt
& \result{85.55}{0.95} & \result{79.87}{0.38}
& \result{81.04}{1.61} & \result{70.59}{1.31}
& \result{71.60}{0.34} & \result{66.50}{0.97}
& \result{52.53}{1.24} & \result{42.33}{0.77}
& \result{75.66}{1.27} & \result{66.71}{0.68}
& \result{86.70}{2.66} & \result{75.86}{3.82}
\\

CODAPrompt
& \result{87.36}{0.76} & \result{81.63}{0.57}
& \result{82.09}{1.74} & \result{70.56}{1.18}
& \result{78.55}{0.61} & \result{73.20}{0.47}
& \result{51.67}{0.49} & \result{41.12}{0.54}
& \result{77.89}{0.93} & \result{69.30}{0.55}
& \result{85.79}{3.42} & \result{74.90}{2.48}
\\

APER-Adapter
& \result{89.54}{0.88} & \result{84.90}{0.31}
& \result{91.63}{1.12} & \result{86.77}{0.03}
& \result{75.76}{0.78} & \result{68.24}{0.50}
& \result{59.05}{0.91} & \result{48.94}{0.14}
& \result{80.48}{1.17} & \result{74.07}{0.29}
& \result{90.60}{0.96} & \result{84.34}{0.20}
\\

EASE
& \result{90.82}{0.67} & \result{85.97}{0.70}
& \result{90.17}{1.21} & \result{84.47}{1.06}
& \result{82.30}{0.55} & \result{76.42}{0.21}
& \result{57.29}{0.86} & \result{45.11}{1.62}
& \result{76.04}{1.59} & \result{68.31}{0.54}
& \result{89.78}{1.21} & \result{82.94}{0.78}
\\

MOS
& \result{93.15}{0.63} & \result{89.73}{0.23}
& \result{93.54}{0.52} & \bestresult{90.15}{0.18}
& \result{82.98}{0.69} & \result{78.59}{0.35}
& \result{61.64}{1.80} & \result{53.89}{1.11}
& \result{85.86}{1.00} & \result{80.08}{0.09}
& \result{95.44}{0.36} & \result{91.78}{0.55}
\\

\midrule
\rowcolor{toscarow}

TOSCA (ours)
& \bestresult{94.55}{0.67} & \bestresult{91.50}{0.32}
& \bestresult{93.63}{0.63} & \result{90.13}{0.19}
& \bestresult{85.94}{0.30} & \bestresult{81.45}{0.21}
& \bestresult{67.68}{1.88} & \bestresult{59.88}{1.00}
& \bestresult{86.56}{0.77} & \bestresult{80.46}{0.30}
& \bestresult{96.32}{0.45} & \bestresult{93.28}{0.06}
\\

\bottomrule
\end{tabular}%
}

\end{table*}

\begin{table*}[htbp]
\centering
\caption{Efficiency comparison on CIFAR B0 Inc5 with ViT-B/16 pretrained on ImageNet-21K, reported as mean $\pm$ standard deviation over five seeds. Training FLOPs are measured directly with \texttt{torch.profiler} over real forward and backward passes rather than estimated analytically. Columns without a standard deviation are deterministic given the architecture and task split. The lowest cost in each column is shown in bold.}
\label{tab:efficiency}
\resizebox{\textwidth}{!}{%
\begin{tabular}{lccccccc}
\toprule
Method & Train Time (min) & Eval. (ms/smp) & Inference (GFLOPs) & Train (GFLOPs) & Peak Train Mem. (MB) & Params (M) & Trainable Params (M) \\
\midrule
SimpleCIL & $\mathbf{2.17 \pm 0.03}$ & $\mathbf{2.42 \pm 0.01}$ & $\mathbf{35.14}$ & $\mathbf{1.76 \times 10^{6}}$ & $\mathbf{1384.09}$ & $\mathbf{85.88}$ & $85.88$ \\
RanPAC & $12.28 \pm 0.11$ & $2.65 \pm 0.01$ & $35.62$ & $3.98 \times 10^{7}$ & $4415.97$ & $87.99$ & $2.19$ \\
L2P & $46.43 \pm 0.06$ & $5.65 \pm 0.00$ & $74.91$ & $2.91 \times 10^{7}$ & $2362.72$ & $171.82$ & $\mathbf{0.12}$ \\
DualPrompt & $42.13 \pm 0.05$ & $5.13 \pm 0.01$ & $70.29$ & $2.65 \times 10^{7}$ & $2832.75$ & $172.00$ & $0.33$ \\
CODAPrompt & $185.29 \pm 0.32$ & $4.62 \pm 0.00$ & $70.29$ & $1.36 \times 10^{8}$ & $18726.53$ & $89.72$ & $89.72$ \\
APER-Adapter & $9.28 \pm 0.08$ & $4.91 \pm 0.01$ & $70.74$ & $7.46 \times 10^{7}$ & $4415.97$ & $172.94$ & $87.14$ \\
EASE & $113.95 \pm 2.16$ & $49.03 \pm 0.02$ & $707.69$ & $4.44 \times 10^{8}$ & $4535.86$ & $88.53$ & $2.73$ \\
MOS & $100.19 \pm 2.90$ & $52.80 \pm 0.17$ & $740.37$ & $1.40 \times 10^{8}$ & $4650.50$ & $92.34$ & $6.47$ \\
\midrule
\rowcolor{toscarow}
TOSCA (ours) & $50.72 \pm 0.15$ & $3.62 \pm 0.01$ & $35.65$ & $7.39 \times 10^{7}$ & $25351.67$ & $101.73$ & $0.23$ \\
\bottomrule
\end{tabular}%
}
\end{table*}

\bibliographystyle{plainnat}
\bibliography{references}
\end{document}